% Source pp. 264-272 (Lectures19-25.pdf).
\section{Energy-Momentum Operators}

The expression \eqref{eq:Pcomm} can be now promoted to the definition of the
operators $\hat P^\mu$ acting in the space $\Ham$. Specifically, for any
$x_1,\dots,x_N,y_1,\dots,y_M$ in $\R^d_-$ define
\begin{equation}\label{eq:Pmatrix}
\bigl\langle\varphi(y_1)\cdots\varphi(y_M)\big|\hat P^\mu\big|\varphi(x_1)\cdots\varphi(x_N)\bigr\rangle
=\bigl\langle\varphi(\bar y_1)\cdots\varphi(\bar y_M)\,P^\mu\,\varphi(x_1)\cdots\varphi(x_N)\bigr\rangle
\end{equation}
where the correlation function in the r.h.s.\ involves the integral of the
energy-momentum tensor over the equal-time slice $\tau=0$,
\begin{equation}\label{eq:Pdef}
P^\mu(\tau):=-i\int_\Sigma \dd\Sigma_\nu(x)\,T^{\mu\nu}(x)=-i\int_{\Sigma(\tau=0)}\dd^{d-1}\xx\,T^{\mu\tau}(\xx,\tau=0)\,,
\end{equation}
It is straightforward to check, using the definition of the inner product, that the
``spatial'' components $\PPv=(P^1,\dots,P^{d-1})$ are Hermitian, while $P^\tau$ is
anti-Hermitian. We therefore define $d$ Hermitian operators $(\hat{\PPv},\hat H)$ as
\begin{align}
\hat P^k&:=-i\int_{\Sigma(\tau=0)}\dd^{d-1}\xx\,T^{k\tau}(\xx,\tau)\,,\label{eq:Pk}\\
\hat H&=-\int_{\Sigma(\tau=0)}\dd^{d-1}\xx\,T^{\tau\tau}(\xx,\tau)\,.\label{eq:Hdef}
\end{align}
Obviously, $\hat H$ coincides with $\hat P^0:=-i\hat P^\tau$, the real-time
component of $\hat P^\mu$.

A remark is in order. If we are dealing with the field theory in the infinite
space-time $\R^d$, the matrix elements of these operators diverge because of the
contribution of generally nonzero expectation value
\[
\bigl\langle T^{\mu\nu}(x)\bigr\rangle=g^{\mu\nu}\,\varepsilon
\]
where $g^{\mu\nu}=\delta^{\mu\nu}$ is the Euclidean metric tensor, and
$\varepsilon$ is the vacuum energy density. It is convenient to get rid of this
divergence by defining reduced energy-momentum tensor
\[
T^{\mu\nu}(x)\to T^{\mu\nu}(x)-g^{\mu\nu}\,\varepsilon
\]
which has zero expectation value. This can be interpreted as adding the
``cosmological term'' to the action
\[
\mathcal A\to\mathcal A+\varepsilon\int \dd^dx\,\sqrt{g(x)}
\]
For the fixed metric background such term has no effect on the content of the
theory, but brings in the above subtraction when $T^{\mu\nu}(x)$ is defined as the
response to the variation of the background metric,
\[
\delta\mathcal A=-\frac12\int \dd^dx\,\sqrt g\,T^{\mu\nu}(x)\,\delta g_{\mu\nu}(x)\,.
\]
In what follows we adopt this convention, assuming that $T^{\mu\nu}(x)$ is chosen in
such a way that
\begin{equation}\label{eq:Tvev0}
\bigl\langle T^{\mu\nu}(x)\bigr\rangle=0\,.
\end{equation}

With this convention, it is easy to prove that the vacuum $\ket0$ is the eigenstate
of $\hat P^\mu$ with the eigenvalue 0, i.e.
\begin{equation}\label{eq:vacP}
\hat H\ket0=0\,,\qquad\hat{\PPv}\ket0=0\,.
\end{equation}
Indeed, consider the element
\begin{equation}\label{eq:vacelement}
\bigl\langle O_1(x_1)\cdots O_N(x_N)\big|\hat P^\mu\big|0\bigr\rangle=\bigl\langle O_1(\bar x_1)\cdots O_N(\bar x_N)\,P^\mu\bigr\rangle
\end{equation}
where $P^\mu$ is the surface integral \eqref{eq:Pdef}. As all insertions
$\bar x_1,\dots,\bar x_N$ are in $\R^d_+$, the hypersurface $\Sigma$ involved in the
definition \eqref{eq:Pmatrix} can be deformed away to the domain $\tau\to-\infty$,
where the integral vanishes. As the result, all matrix elements
\eqref{eq:vacelement} are equal to zero. Eq.~\eqref{eq:vacP} follows.

Similar calculation shows that the operators $P_\mu$ act on the states
$\ket{O_1(x_1)\cdots O_N(x_N)}$ as follows
\[
\hat P_\mu\ket{O_1(x_1)\cdots O_N(x_N)}=i\left[\sum_{k=1}^N\frac\partial{\partial x_k^\mu}\right]\ket{O_1(x_1)\cdots O_N(x_N)}
\]
Indeed, this time deforming $\Sigma$ to the $\tau\to-\infty$ domain leaves behind
integrals around the insertion points $x_k$, which, according to
\[
\partial^\mu O(x)=\int_{\partial D_x}\dd\Sigma_\nu(y)\,T^{\mu\nu}(y)\,O(x)
\]
leaves behind the derivatives in the r.h.s. As usual, the infinitesimal
translations generated by $\hat P_\mu$ can be exponentiated to finite translations.
In terms of the operators \eqref{eq:Pk}, \eqref{eq:Hdef} we have
\begin{equation}\label{eq:finitetr}
\begin{aligned}
e^{-i\hat{\PPv}\cdot\mathbf X}\ket{O_1(\xx_1,\tau_1)\cdots O_N(\xx_N,\tau_N)}&=\ket{O_1(\xx_1+\mathbf X,\tau_1)\cdots O_N(\xx_N+\mathbf X,\tau_N)}\\
e^{-\hat HL}\ket{O_1(\xx_1,\tau_1)\cdots O_N(\xx_N,\tau_N)}&=\ket{O_1(\xx_1,\tau_1-L)\cdots O_N(\xx_N,\tau_N-L)}
\end{aligned}
\end{equation}

\subsection*{Stationary States. Energy Positivity}

The operators $\hat P^\mu$ commute, and thus can be simultaneously diagonalized in
$\Ham$. It is convenient to discuss the spectrum of the Hermitian operators
$\hat P^\mu$, $\mu=0,\dots,d-1$ with $\hat P^0:=-i\hat P^\tau=\hat H$, which, of
course, coincide with the conventional Energy-Momentum operators in the Minkowski
space-time. Let $\ket\alpha$ be an eigenstate
\[
\hat P^\mu\ket\alpha=p^\mu_\alpha\ket\alpha\,.
\]
We will make few assumptions about the spectrum of these operators. One is about
the spectrum of $\hat H$. We assume that the eigenvalues $E_\alpha$ of $\hat H$
satisfy $E_\alpha\ge0$, with the only zero eigenvalue $E_0=0$ being associated with
the vacuum state, Eq.~\eqref{eq:vacP}. This assumption implies
\begin{equation}\label{eq:projector}
\lim_{L\to\infty}e^{-L\hat H}=\ket0\bra0
\end{equation}
There are several simple consequences from this condition alone. One is the
so-called cluster factorization property of the correlation functions, which states
that if a collection of $N$ points consists of two groups with very large separation
between the groups, correlations between the fields in the different groups tend to
zero, that is
\begin{multline}\label{eq:cluster}
\lim_{L\to\infty}\bigl\langle O_1(x_1)\cdots O_M(x_M)\,O_{M+1}(x_{M+1}+L)\cdots O_N(x_N+L)\bigr\rangle\\
=\bigl\langle O_1(x_1)\cdots O_M(x_M)\bigr\rangle\,\bigl\langle O_{M+1}(x_{M+1})\cdots O_N(x_N)\bigr\rangle
\end{multline}
To prove, take the direction of $L$ to be the ``time'', then the l.h.s.\ can be
represented as
\[
\bigl\langle O_1(x_1)\cdots O_M(x_M)\big|e^{-L\hat H}\big|O_{M+1}(x_{M+1})\cdots O_N(x_N)\bigr\rangle
\]
and \eqref{eq:cluster} follows from \eqref{eq:projector}.

Less obvious consequence concerns analytic properties of the correlation functions
\[
\bigl\langle O_1(\xx_1,\tau_1)\cdots O_N(\xx_N,\tau_N)\bigr\rangle
\]
at complex values of $\tau_1,\dots,\tau_N$. We have briefly studied this question at
the beginning of the course, for the case of the Klein-Gordon propagator
(Section~4). In fact, our arguments there depend only on the condition of the
positivity of the energy spectrum, which was explicit in the KG theory. Therefore,
the analysis can be repeated word for word in general case. Let me restate the
result for the two point function
\[
\bigl\langle O_1(\xx,\tau)\,O_2(\mathbf0,0)\bigr\rangle\,.
\]
In the complex plane of $t=-i\tau$ we have the picture of
figure~\ref{fig:tcontours2}.
\begin{figure}[ht]
\centering\figfile{tcontours2}
\caption{The two-point function in the complex $t$-plane. Above the light-cone
cuts (along $C_+$) it gives $\bra0\hat O_2(\mathbf0,0)\,\hat O_1(\xx,t)\ket0$,
below them (along $C_-$) it gives $\bra0\hat O_1(\xx,t)\,\hat O_2(\mathbf0,0)\ket0$.}
\label{fig:tcontours2}
\end{figure}
Here
\[
\hat O(\xx,t)=e^{-i\PPv\cdot\xx+iHt}\,\hat O(\mathbf0,0)\,e^{i\PPv\cdot\xx-iHt}
\]
are conventional Heisenberg field operators. The discontinuity across the branch
cuts represents the commutators which lie within the future and past light cones
(``local commutativity'').

\subsection*{Spectrum of Energy-Momentum}

Some properties of the spectrum of the operator $\hat P^\mu=(\hat H,\hat{\PPv})$ can
be established on general grounds. Let $\mathcal P$ be the spectrum of this
operator, i.e.\ the set of its eigenvalues
\[
p^\mu\in\mathcal P\ \Rightarrow\ \exists\ \ket\alpha\in\Ham\ \text{such that}\ p^\mu=p^\mu_\alpha:\quad\hat P^\mu\ket\alpha=p^\mu_\alpha\ket\alpha\,.
\]
Lorentz invariance suggests that
\[
p^\mu\in\mathcal P\ \Rightarrow\ \Lambda^\mu{}_\nu\,p^\nu\in\mathcal P\,,
\]
where $\Lambda$ is a Lorentz transformation. This means that any $p^\mu\in\mathcal P$
such that $p^\mu p_\mu=-m^2$ enters $\mathcal P$ together with the entire hyperboloid
\[
\mathcal P_{m^2}=\bigl\{p^\mu:\ -p^\mu p_\mu=(p^0)^2-\pp^2=m^2\bigr\}\,.
\]
It follows then from the positivity of the energy that $\mathcal P$ lies entirely
within the future light cone $p^\mu p_\mu\le0$, $p^0\ge0$.

Another important property of $\mathcal P$ is its additivity, i.e.
\[
p^\mu_\alpha,\,p^\mu_\beta\in\mathcal P\ \Rightarrow\ p^\mu_\alpha+p^\mu_\beta\in\mathcal P\,.
\]
Intuitively, if we think of the state $\ket\alpha$ carrying the momentum
$p^\mu_\alpha$ as a local excitation over the ground state, it is possible to add
another excitation, with the momentum $p^\mu_\beta$ arbitrarily far away from the
first one. At large separations the excitations do not interact, and their energies
and momenta add up. Formal derivation using the cluster property will be considered
later.

The additivity implies that if $\mathcal P$ contains $\mathcal P_{m^2}$, it also
contains all $p^\mu$ of the form
\[
p^\mu=p_1^\mu+\dots+p_N^\mu\,;\qquad p_i^\mu\in\mathcal P_{m^2}\,.
\]
The totality of such vectors $p^\mu$ will be denoted $\bar{\mathcal P}_{m^2}$,
\[
\bar{\mathcal P}_{m^2}=\bigcup_N\ \bigcup_{p_1,\dots,p_N}\bigl\{p^\mu=p_1^\mu+\dots+p_N^\mu;\ p_i^\mu\in\mathcal P_{m^2}\bigr\}\,.
\]
The set $\bar{\mathcal P}_{m^2}$ looks like figure~\ref{fig:spectrum} of
Section~4.

Note that $\bar{\mathcal P}_{m^2}$ is exactly the spectrum of $\hat P^\mu$ in the
Klein-Gordon field theory, where we have found that
\[
\hat P^\mu=\int\frac{\dd^3\pp}{(2\pi)^3}\,p^\mu\,a^\dagger_\pp\,a_\pp\,,
\]
where $p^\mu=(\omega_\pp,\pp)$, $\omega_\pp=\sqrt{\pp^2+m^2}$, and $a_\pp$,
$a^\dagger_\pp$ are particle creation and annihilation operators. There these
states are interpreted as $N$-particle states of free bosons.

The spectrum $\bar{\mathcal P}_{m^2}$ in some sense is the minimal Lorentz-invariant
set compatible with the above additivity condition. If this spectrum appears in an
interacting QFT we say that the theory contains just one massive particle of the
mass $m$. The states associated with the hyperboloid $\mathcal P_{m^2}$ are the
one-particle states, and the states with $-p^\mu p_\mu\ge4m^2$ are interpreted as the
scattering states of two or more interacting particles.

In general case the spectrum can be more complicated (even in a scalar QFT like we
consider here). It can involve more than one kind of stable particles, and hence
would contain several hyperboloids, say $\mathcal P_{m_1^2}$ and
$\mathcal P_{m_2^2}$, and correspondingly its continuous part would look more
complicated, as in figure~\ref{fig:spectrum2}.
\begin{figure}[ht]
\centering\figfile{spectrum2}
\caption{Energy-momentum spectrum of a theory with two stable particles of masses
$m_1<m_2$: two mass hyperboloids, and a continuum of scattering states starting at
$\min(2m_1,m_1+m_2)=2m_1$, with further thresholds at $m_1+m_2$, $2m_2$, $3m_1$,
\ldots{} inside it.}
\label{fig:spectrum2}
\end{figure}

In any case, if we form the operator
\[
-\hat P^\mu\hat P_\mu=\hat H^2-\hat{\PPv}^2\,,
\]
the discrete components of its spectrum are associated with the one-particle
states, and the continuous spectrum incorporates all the multi-particle scattering
states.

``Asymptotic completeness'': $\mathcal P=\cup_a\bar{\mathcal P}_{m_a^2}$.

Lorentz invariance of the spectrum has certain implications for the general
structure of the correlation functions in QFT. As the simplest example, consider the
two-point correlation function
\[
\bigl\langle\varphi(\xx,\tau)\,\varphi(\mathbf0,0)\bigr\rangle=\bra0T_\tau\bigl(\hat\varphi(\xx,\tau)\,\hat\varphi(\mathbf0,0)\bigr)\ket0\,,
\]
where $(\xx,\tau)$ are the Euclidean space-time coordinates. Assuming $\tau>0$ we can
write it as
\[
\sum_\alpha\bra0\hat\varphi(\xx,\tau)\ket\alpha\bra\alpha\hat\varphi(\mathbf0,0)\ket0\,.
\]
Using the fact that
\[
\hat\varphi(\xx,\tau)=e^{\hat H\tau-i\hat{\PPv}\cdot\xx}\,\hat\varphi(\mathbf0,0)\,e^{-\hat H\tau+i\hat{\PPv}\cdot\xx}\,.
\]
we can write the above sum as
\[
\sum_\alpha\bra0\hat\varphi(\mathbf0,0)\ket\alpha\bra\alpha\hat\varphi(\mathbf0,0)\ket0\,e^{-E_\alpha\tau+i\PPv_\alpha\cdot\xx}\,,
\]
where $E_\alpha$, $\PPv_\alpha$ are the energy and momentum of the state
$\ket\alpha$. In fact, we have continuously many states, so that
\[
\sum_\alpha\to\int \dd\mu(\alpha)\,.
\]
where $\alpha$ denotes complete set of quantum numbers characterizing the state.
Define
\begin{equation}\label{eq:rhodef}
\rho(q)=\int \dd\mu(\alpha)\,\bigl|\bra0\hat\varphi(0)\ket\alpha\bigr|^2\,(2\pi)^4\,\delta^{(4)}(q^\mu-p^\mu_\alpha)
\end{equation}
Then, assuming $\tau>0$, we have
\begin{equation}\label{eq:spectralrep}
\bigl\langle\varphi(\xx,\tau)\,\varphi(\mathbf0,0)\bigr\rangle=\int \dd\mu(\alpha)\,\bra0\hat\varphi(\xx,\tau)\,\hat\varphi(\mathbf0,0)\ket0
=\int\frac{\dd^4q}{(2\pi)^4}\,\rho(q)\,e^{-q^0\tau+i\qq\cdot\xx}
\end{equation}
Evidently,
\[
\rho(q)=0\quad\text{if}\quad q^\mu\notin\mathcal P\,,\qquad\rho(q)\ge0\quad\text{if}\quad q^\mu\in\mathcal P\,,
\]
so that $\rho(q)\ge0$ everywhere.

$\rho(q)$ vanishes if $q^0<0$ (Hamiltonian is assumed to be positive definite), and
at $q^0>0$ it depends on $-q^2=(q^0)^2-\qq^2$ (Lorentz invariance); I will write
\[
\rho=\theta(q^0)\,\rho(-q^2)\,.
\]
Then, for $\tau>0$ we have
\begin{align*}
\bigl\langle\varphi(\xx,\tau)\,\varphi(\mathbf0,0)\bigr\rangle
&=\int_0^\infty \dd\mu^2\,\rho(\mu^2)\int\frac{\dd^4q}{(2\pi)^4}\,\delta\bigl((q^0)^2-\qq^2-\mu^2\bigr)\,\theta(q^0)\,e^{-q^0\tau+i\qq\cdot\xx}\\
&=\int_0^\infty\frac{\dd\mu^2}{2\pi}\,\rho(\mu^2)\int\frac{\dd^3\qq}{(2\pi)^3}\,\frac1{2\omega_\qq}\,e^{-\omega_\qq\tau+i\qq\cdot\xx}\,,
\end{align*}
where
\[
\omega_\qq=\sqrt{\qq^2+\mu^2}\,.
\]
Proceeding similarly, the case $\tau<0$ we obtain the same expression with
\[
e^{-\omega_\qq\tau}\to e^{\omega_\qq\tau}
\]
Both cases $\tau>0$ and $\tau<0$ can be put together by using
\[
\frac1{2\omega_q}\,e^{-\omega_q|\tau|}=\int\frac{\dd q_4}{2\pi}\,\frac{e^{iq_4\tau}}{q_4^2+\qq^2+\mu^2}
\]
Finally
\begin{equation}\label{eq:KL}
\bigl\langle\varphi(\xx,\tau)\,\varphi(\mathbf0,0)\bigr\rangle=\int_0^\infty\frac{\dd\mu^2}{2\pi}\,\rho(\mu^2)\,D(|x|,\mu^2)\,,
\end{equation}
where $|x|=\sqrt{\tau^2+\xx^2}$, and $D(|x|,\mu^2)$ is the Klein-Gordon propagator
with the mass $\mu$,
\[
D(|x|,\mu^2)=\int\frac{\dd^4k}{(2\pi)^4}\,\frac{e^{i\kk\cdot\xx+ik_4\tau}}{k_4^2+\kk^2+\mu^2}\,.
\]
Eq.~\eqref{eq:KL} is known as the \emph{K\"all\'en-Lehmann representation}.

For $\bar{\mathcal P}_{m^2}$
\[
\rho(\mu^2)=(2\pi)\,Z\,\delta(\mu^2-m^2)+\tilde\rho(\mu^2)
\]
where $\tilde\rho(\mu^2)=0$ at $\mu^2<4m^2$,
\begin{equation}\label{eq:Zres}
Z^{1/2}=\bra0\hat\varphi(0)\ket\pp
\end{equation}
\begin{equation}\label{eq:KLmom}
\tilde W^{(2)}(p)=\frac Z{p^2+m^2}+\int_{4m^2}^\infty\frac{\tilde\rho(\mu^2)}{p^2+\mu^2}\,\frac{\dd\mu^2}{2\pi}
\end{equation}
which is another form of the K\"all\'en-Lehmann representation.

Note the pole at $p^2=-m^2$ explicit in this formula, appearing entirely from the
contributions of the one-particle states in the intermediate-state decomposition.
The residue $Z$ is positive (as it is the square of the absolute value of the matrix
element \eqref{eq:Zres}), but its value depends on the normalization of the field
$\varphi$. That normalization was subject to one of the normalization conditions in
our calculations in the perturbation theory. Thus, under the normalization
condition
\begin{equation}\label{eq:SIrecall}
\tilde\Gamma^{(2)}(p^2)=p^2+m^2+O\bigl((p^2+m^2)^2\bigr)
\end{equation}
of the scheme SI in the $\varphi^4$ theory (see Eqs.~\eqref{eq:SIa},
\eqref{eq:SIW}) we just choose $Z=1$. Let us also note that $\tilde\rho(\mu^2)$ is
non-negative at all $\mu^2$, as the consequence of the positivity of the norms in
the Hilbert space, see Eq.~\eqref{eq:rhodef}.

At $|x|\to\infty$ we have
\[
D(|x|,\mu^2)\simeq\frac{\mu^{1/2}}{|x|^{3/2}}\,e^{-\mu|x|}\,,
\]
and the form \eqref{eq:KL} shows that at large $|x|$ the asymptotic of
$\langle\varphi(x)\varphi(0)\rangle$ is always dominated by the one-particle
contribution, provided the matrix element \eqref{eq:Zres} is not zero.
